\chapter{Theoretischer Hintergrund}
\label{chap:theoretischer-hintergrund}

\section{Grundlagen von E-Commerce und Online-Marktplätzen}
\label{sec:grundlagen-ecommerce}

\subsection{Begriffsdefinitionen: E-Commerce, Marktplatz und Multi-Rollen-Modell}
\label{subsec:begriffe}

E-Commerce bezeichnet den elektronischen Handel mit Waren und
Dienstleistungen über digitale Netzwerke, insbesondere das Internet.
Er umfasst sämtliche onlinebasierten Geschäftstransaktionen zwischen
Unternehmen und Konsumenten (\acfirst{B2C}) sowie zwischen Unternehmen
untereinander (\acfirst{B2B}), einschließlich der digitalen Vermittlung
von Dienstleistungen und der Abwicklung zugehöriger Geschäftsprozesse
\cite{turban_ecommerce}.

Ein Online-Marktplatz ist eine spezielle Form des E-Commerce, bei
der eine Plattform als Intermediär zwischen mehreren Anbietern und
Nachfragern auftritt. Im Unterschied zum klassischen Webshop eines
einzelnen Händlers bündelt ein Marktplatz das Angebot verschiedener
Anbieter und stellt die technische sowie organisatorische
Infrastruktur für deren Interaktion bereit
\cite{parker_platform_revolution}. Die Plattform übernimmt dabei
Funktionen wie Authentifizierung, Zahlungsabwicklung,
Bewertungssysteme und Zugriffskontrolle
\cite{tiwana_platform_ecosystems}.

Ein Multi-Rollen-Modell erweitert dieses Konzept, indem neben
Käufern und Verkäufern weitere Akteursgruppen wie Dienstleister
und Administratoren in das System integriert werden.
Jede Rolle besitzt spezifische Rechte, Verantwortlichkeiten und
Interaktionsmöglichkeiten innerhalb der Plattform. Dies erfordert
eine rollenbasierte Zugriffskontrolle sowie klar definierte
Workflows, um eine sichere und konsistente Interaktion zwischen
den Akteuren zu gewährleisten \cite{parker_platform_revolution,
tiwana_platform_ecosystems}.

\subsection{Marktplatztypen und Geschäftsmodelle}
\label{subsec:marktplatztypen}

Digitale Marktplätze lassen sich nach verschiedenen Kriterien
klassifizieren. Eine verbreitete Unterscheidung erfolgt anhand
der Anzahl der beteiligten Marktseiten. Einseitige Marktplätze
verbinden einen Anbieter direkt mit Endkunden, während zweiseitige
Marktplätze zwei voneinander abhängige Nutzergruppen zusammenführen.
Multi-Sided Plattformen verbinden darüber hinaus mehr als zwei
Gruppen miteinander und ermöglichen komplexere Interaktionsmuster
\cite{heinemann_online_handel}.

Eine weitere Klassifikation unterscheidet zwischen vertikalen
Marktplätzen, die sich auf eine bestimmte Branche spezialisieren,
und horizontalen Marktplätzen, die branchenübergreifend verschiedene
Produkte und Dienstleistungen anbieten. \textit{Bazario} folgt
dem horizontalen Modell, da es sowohl Produkte als auch
Dienstleistungen verschiedener Anbieter in einem einheitlichen
System zusammenführt \cite{taeuscher_platform_models}.

Hinsichtlich der Geschäftsmodelle lassen sich provisionsbasierte
Modelle, bei denen die Plattform einen prozentualen Anteil am
Transaktionswert einbehält, feste Transaktionsgebühren sowie
hybride Ansätze unterscheiden, die beide Varianten kombinieren.
\textit{Bazario} setzt auf ein hybrides Modell: Neben einer
festen, unabhängig vom Umsatz erhobenen Grundprovision fällt
eine zusätzliche, transaktionsabhängige Provision an, deren
Höhe sich nach dem über die Plattform abgewickelten
Bestellvolumen richtet. Beide Anteile werden dem Anbieter bei
jeder Transaktion abgezogen. Die konkrete Konfiguration der
Provisionssätze erfolgt über das Backend und ist im Rahmen
dieser Arbeit nicht Gegenstand des Frontends.
Tabelle~\ref{tab:abrechnungsmodelle} stellt die drei verbreiteten 
Abrechnungsmodelle digitaler Marktplätze vergleichend gegenüber. 
Die Merkmale wurden auf Basis von \cite{taeuscher_platform_models} 
und \cite{heinemann_online_handel} erarbeitet und zusammengefasst.

\begin{table}[H]
\centering
\caption{Vergleich digitaler Abrechnungsmodelle}
\label{tab:abrechnungsmodelle}
\begin{tabular}{|l|p{5.5cm}|p{5.5cm}|}
\hline
\textbf{Modell} & \textbf{Vorteile} & \textbf{Nachteile} \\
\hline
Provisionsbasiert &
Automatisch umsatzabhängig, fair für kleine Anbieter &
Geringere Planbarkeit der Einnahmen \\
\hline
Feste Transaktionsgebühr &
Planbare Einnahmen, einfache Abrechnung &
Hohe Belastung bei geringwertigen Transaktionen \\
\hline
Hybrid &
Ausgewogene Balance zwischen Stabilität und Flexibilität &
Höhere Komplexität in der Abrechnung \\
\hline
\end{tabular}
\end{table}


\subsection{Grundlagen der Online-Zahlung}
\label{subsec:online-zahlung}

Die Zahlungsabwicklung ist ein zentraler Bestandteil jeder
E-Commerce-Plattform. Moderne Zahlungssysteme basieren häufig
auf einer intentionsbasierten Architektur, bei der eine Zahlung
zunächst als sogenannter Payment Intent vorbereitet und autorisiert
wird, bevor sie final ausgeführt wird. Dieser Ansatz ermöglicht
eine flexible Kopplung des Zahlungsprozesses an Geschäftsprozesse
wie Reservierungen, Buchungen oder Bestellbestätigungen
\cite{stripe_security_guide}.

Für die Integration im Frontend stellt Stripe mit Stripe Elements
eine Bibliothek vorgefertigter, \acs{PCI-DSS}-konformer \acs{UI}-Komponenten
bereit, die eine sichere Erfassung von Zahlungsdaten direkt im
Browser ermöglichen, ohne dass sensible Kartendaten den eigenen
Server passieren \cite{stripe_elements_docs}.

Neben dieser eingebetteten Integration bietet Stripe mit
Stripe Checkout eine alternative, gehostete Zahlungsseite an:
Das Frontend leitet den Nutzer dabei lediglich auf eine von
Stripe bereitgestellte, vorkonfigurierte Zahlungsseite weiter
und verarbeitet nach Abschluss der Zahlung die Rückleitung
auf eine Erfolgs- oder Abbruchseite. Im Gegensatz zur
eingebetteten Variante entfällt dadurch die Notwendigkeit,
eigene Zahlungsformulare im Frontend zu gestalten und zu
validieren \cite{stripe_checkout_docs}.

Der \acfirst{PCI-DSS} ist ein
Sicherheitsstandard, der Anforderungen an die sichere Verarbeitung,
Speicherung und Übertragung von Kartendaten definiert. Sowohl bei Stripe Elements als auch bei Stripe Checkout
werden sensible Zahlungsdaten ausschließlich auf den Servern
von Stripe verarbeitet, wodurch der \acs{PCI-DSS}-Scope des
Betreibers erheblich reduziert wird
\cite{stripe_security_guide}.

\subsection{Wettbewerbsanalyse bestehender Marktplatzplattformen}
\label{subsec:wettbewerbsanalyse}

Im deutschen E-Commerce-Markt dominieren sowohl internationale
als auch lokale Plattformen. Marktführer wie Amazon verbinden
Käufer, Verkäufer, Logistikpartner und Administratoren in einem
Multi-Sided-Modell und setzen dabei auf starke Kundenbindung
sowie eine breite Produktauswahl \cite{parker_platform_revolution}.
Spezialisierte Plattformen wie Booking.com zeigen, dass auch
vertikale Marktplätze mit Fokus auf Dienstleistungsbuchungen
hohe Marktanteile erzielen können: Der Mutterkonzern Booking
Holdings erzielte im Geschäftsjahr 2024 einen Gesamtumsatz von
23,7 Milliarden US-Dollar bei rund 1,1 Milliarden vermittelten
Übernachtungen \cite{booking_holdings_q4_2024}.

Tabelle~\ref{tab:wettbewerb} vergleicht ausgewählte bestehende
Plattformen hinsichtlich ihrer Kernfunktionen und zeigt, wie
sich \textit{Bazario} von diesen abgrenzt. Die Einordnung
orientiert sich an den Klassifikationskriterien von
\cite{taeuscher_platform_models} und \cite{heinemann_online_handel}.

\begin{table}[h]
\centering
\caption{Vergleich bestehender Marktplatzplattformen}
\label{tab:wettbewerb}
\begin{tabular}{|p{2.8cm}|p{2.5cm}|p{2.5cm}|p{2.5cm}|p{2.5cm}|}
\hline
\textbf{Plattform} & \textbf{Produkte} &
\textbf{Dienst-leistungen} & \textbf{Multi-Rolle} &
\textbf{Buchungs-system} \\
\hline
Amazon & \checkmark & eingeschränkt & \checkmark & \texttimes \\
\hline
Etsy & \checkmark & \texttimes & eingeschränkt & \texttimes \\
\hline
Booking.com & \texttimes & \checkmark & eingeschränkt &
\checkmark \\
\hline
Fiverr & \texttimes & \checkmark & eingeschränkt & \checkmark \\
\hline
\textit{Bazario} & \checkmark & \checkmark & \checkmark &
\checkmark \\
\hline
\end{tabular}
\end{table}

\section{Frontend-Architektur für webbasierte Anwendungen}
\label{sec:frontend-architektur}

\subsection{Single-Page Applications vs.\ Multi-Page Applications}
\label{subsec:spa-mpa}

Webbasierte Anwendungen lassen sich grundlegend in zwei
Architekturansätze unterteilen: Multi-Page Applications
(\acfirst{MPA}) und Single-Page Applications (\acfirst{SPA}). Bei \acs{MPA}s wird bei jeder
Nutzerinteraktion eine vollständige \acs{HTML}-Seite vom Server geladen
und im Browser neu gerendert. Dieser Ansatz führt bei komplexen
Anwendungen zu spürbaren Ladezeiten und einer fragmentierten
Nutzererfahrung \cite{bhalla_spa_architecture}.

\acs{SPA}s hingegen laden die Anwendung einmalig und aktualisieren
anschließend nur die jeweils relevanten Teile der Oberfläche.
Die Navigation erfolgt clientseitig über JavaScript, wodurch eine
flüssige, appähnliche Nutzererfahrung entsteht
\cite{jonathan_react_spa}. Tabelle~\ref{tab:spa-mpa} stellt
beide Ansätze vergleichend gegenüber.


\begin{table}[h]
\centering
\caption{Vergleich SPA und MPA}
\label{tab:spa-mpa}
\begin{tabular}{|p{3.5cm}|p{4.5cm}|p{4.5cm}|}
\hline
\textbf{Kriterium} & \textbf{SPA} & \textbf{MPA} \\
\hline
Seitennavigation & Clientseitig, ohne Neuladen &
Serverseitig, vollständiges Neuladen \\
\hline
Performance & Schnelle Interaktion nach initialem Laden &
Höhere Ladezeiten bei jedem Seitenwechsel \\
\hline
SEO & Eingeschränkt ohne SSR & Gut geeignet \\
\hline
Komplexität & Höher im Frontend & Höher im Backend \\
\hline
Geeignet für & Interaktionsreiche Anwendungen &
Inhaltsreiche Webseiten \\
\hline
\end{tabular}
\end{table}

Für interaktionsreiche Plattformen wie \textit{Bazario}, bei
denen verschiedene Rollen unterschiedliche Ansichten und Workflows
erfordern, bietet die \acs{SPA}-Architektur entscheidende Vorteile
hinsichtlich Performance und Benutzererfahrung. Da
\textit{Bazario} keinen primären SEO-Fokus verfolgt, wiegen
die Vorteile der \acs{SPA}-Architektur die Einschränkungen
hinsichtlich Suchmaschinenoptimierung deutlich auf.

\subsection{Komponentenbasierte Entwicklung}
\label{subsec:komponentenbasiert}

Ein zentrales Prinzip moderner Frontend-Entwicklung ist die
komponentenbasierte Architektur. Die Benutzeroberfläche wird dabei
in eigenständige, wiederverwendbare Einheiten aufgeteilt, die
jeweils ihre eigene Logik, ihren Zustand und ihre Darstellung
kapseln. Dieser Ansatz fördert die Wartbarkeit und
Wiederverwendbarkeit von Code, da Komponenten unabhängig
voneinander entwickelt, getestet und ausgetauscht werden können
\cite{ijlrp_react_components}.

In der Praxis unterscheidet man zwischen Presentational Components, 
die ausschließlich für die Darstellung zuständig sind, und Container 
Components, die Geschäftslogik und Datenzugriff kapseln 
\cite{abramov_presentational_container}. Diese Trennung hält Komponenten überschaubar und
erleichtert das Testen einzelner Einheiten. Die Qualität
komponentenbasierter Anwendungen hängt wesentlich von einer
sauberen Trennung der Verantwortlichkeiten ab. Häufige
Entwurfsprobleme, sogenannte Code Smells, wie zu große
Komponenten oder starke Kopplungen, beeinträchtigen die
Lesbarkeit und Wartbarkeit des Codes
\cite{ferreira_react_smells}.

\subsection{State Management}
\label{subsec:state-management}

In komplexen \acs{SPA}s mit mehreren Nutzerrollen und
funktionsübergreifenden Daten, wie Authentifizierungsstatus,
Nutzerrolle und Warenkorbinhalt, reicht lokaler
Komponentenzustand häufig nicht aus. Globales State
Management ermöglicht es, Anwendungszustände zentral zu
verwalten und konsistent über alle Komponenten hinweg
verfügbar zu machen \cite{bhalla_spa_architecture}.

Zustand ist eine schlanke, hook-basierte
State-Management-Bibliothek für React-Anwendungen. Ein Store
wird direkt als Hook definiert und kann in jeder Komponente
ohne zusätzliche Provider-Wrapper verwendet werden. Dieses
Konzept reduziert die Komplexität der Zustandsverwaltung
erheblich und ermöglicht eine schnelle sowie wartbare
Implementierung \cite{zustand_docs}. Für \textit{Bazario} verwaltet Zustand zentrale Zustände wie den
Authentifizierungsstatus, die Session des Nutzers einschließlich
der zugewiesenen Rollen sowie den Warenkorbinhalt. Diese Zustände werden dabei nicht in einem
einzigen, gemeinsamen Store zusammengefasst, sondern nach dem
Prinzip der Separation of Concerns auf mehrere, fachlich
abgegrenzte Stores verteilt, von denen jeder ausschließlich
für einen bestimmten Bereich der Anwendung zuständig ist. Das
Prinzip geht auf Dijkstra zurück, der es als Technik zur
Beherrschung von Komplexität beschreibt, bei der verschiedene
Aspekte eines Systems bewusst getrennt und jeweils in sich
konsistent betrachtet werden \cite{dijkstra_separation_of_concerns}.
In modernen React-Anwendungen wird dieses Prinzip unter anderem
auf die Aufteilung des globalen Zustands in fachlich
abgegrenzte, unabhängig wartbare Einheiten übertragen
\cite{cheesecakelabs_react_soc}.
Neben dem clientseitigen Anwendungszustand, der Informationen
wie die aktive Nutzerrolle oder den Warenkorbinhalt umfasst, muss
eine \acs{SPA} auch den serverseitigen Datenbestand verwalten, also
Daten, die über eine Programmierschnittstelle (\acfirst{API}) vom
Backend abgerufen werden. Diese
beiden Arten von Zustand unterscheiden sich grundlegend: Während
clientseitiger Zustand vollständig im Browser entsteht und
kontrolliert wird, ist serverseitiger Zustand lediglich eine
zwischengespeicherte Abbildung von Daten, deren eigentliche Quelle
auf dem Server liegt und die daher veralten kann. TanStack Query
ist eine Bibliothek, die auf die Verwaltung dieses serverseitigen
Zustands spezialisiert ist. Sie übernimmt das Abrufen,
Zwischenspeichern und Synchronisieren von \acs{API}-Daten und stellt den
aufrufenden Komponenten zugleich den Zustand jeder Anfrage bereit,
etwa ob diese noch geladen wird, fehlgeschlagen ist oder
erfolgreich Daten geliefert hat \cite{tanstack_query_docs}. Für
\textit{Bazario} ergibt sich damit eine klare Aufteilung: Der
clientseitige Anwendungszustand und der serverseitige
Datenbestand werden durch getrennte, jeweils spezialisierte
Mechanismen verwaltet, während TanStack Query durchgängig für
Letzteren zuständig ist.

\subsection{Routing in Single-Page Applications}
\label{subsec:routing}

Da \acs{SPA}s keine serverseitigen Seitennavigationen durchführen,
übernimmt clientseitiges Routing die Aufgabe, verschiedene
Ansichten der Anwendung anhand der \acfirst{URL} darzustellen. React
Router ist die etablierte Lösung für Routing in
React-Anwendungen und ermöglicht die deklarative Definition
von Routen sowie die Navigation zwischen Ansichten ohne
Seitenneuladen \cite{reactrouter_docs}.

Für Multi-Rollen-Plattformen wie \textit{Bazario} ist
rollenbasierter Seitenschutz ein wesentliches
Sicherheitsmerkmal. Bestimmte Routen werden dabei nur für
Nutzer mit den entsprechenden Berechtigungen zugänglich
gemacht. Dieses Konzept basiert auf dem Prinzip der
rollenbasierten Zugriffskontrolle (\acfirst{RBAC}), das den Zugriff
auf Ressourcen und Funktionen abhängig von der zugewiesenen
Rolle eines Nutzers steuert \cite{ferraiolo_rbac_arxiv}.
In der Praxis werden solche geschützten Routen als 
Wrapper-Komponenten implementiert, die den 
Authentifizierungsstatus und die Rolle des Nutzers prüfen, 
bevor die eigentliche Ansicht gerendert wird 
\cite{wieruch_protected_routes}.


% =====================================================================
% ERSATZ fuer Abschnitt 2.3 in Kapitel/2_Theoretischer_Hintergrund.tex
%
% Der bisherige Abschnitt "Technologiestack" wird vollstaendig durch
% den folgenden ersetzt. Kein "Fuer Bazario" mehr, keine
% Auswahlentscheidung. Nur Funktionsumfang und Gegenueberstellung.
% Die Entscheidung selbst steht in 3.4.
%
% HINWEIS ZU DEN LABELS
% Das Section-Label \label{sec:technologiestack} bleibt absichtlich
% unveraendert, obwohl der Titel jetzt anders lautet. Es wird aus 3.4
% referenziert. Labels sind intern und fuer den Leser unsichtbar;
% ein Umbenennen braechte keinen Vorteil, aber das Risiko toter
% Verweise. Dasselbe gilt fuer subsec:react, subsec:rest-api und
% subsec:ui-bibliotheken, die aus Kapitel 3 und 4 referenziert werden.
%
% Umfang: vorher ca. 3,1 Seiten, jetzt ca. 5 Seiten.
% =====================================================================

\section{Frameworks und Kommunikationsansätze im Frontend}
\label{sec:technologiestack}

% ---------------------------------------------------------------------
\subsection{Aufgaben und Funktionsumfang von Frontend-Frameworks}
\label{subsec:framework-funktionen}

Der Begriff Framework wird im Frontend-Bereich uneinheitlich
verwendet. Gemeinsam ist den gebräuchlichen Lösungen, dass sie die
Erzeugung und Aktualisierung der Benutzeroberfläche von der direkten
Manipulation des Document Object Model (\acfirst{DOM}) entkoppeln.
Anstatt einzelne Änderungen am \acs{DOM} zu programmieren, wird der
gewünschte Zustand der Oberfläche in Abhängigkeit von Daten
beschrieben; die notwendigen Änderungen ermittelt und vollzieht das
Framework selbst. Dieses Vorgehen wird als deklaratives Rendering
bezeichnet und bildet die gemeinsame Grundlage aller im Folgenden
betrachteten Ansätze \cite{comparative_spa_frameworks}.

Erheblich unterscheiden sich die Lösungen jedoch im Funktionsumfang.
Umfassende Frameworks liefern neben dem in
Abschnitt~\ref{subsec:komponentenbasiert} beschriebenen
Komponentenmodell auch Routing, Formularverarbeitung,
\acs{HTTP}-Kommunikation und Werkzeuge zur Abhängigkeitsverwaltung
als integrale Bestandteile aus und geben deren Verwendung weitgehend
vor. Bibliotheken decken dagegen ausschließlich die Darstellungsschicht
ab und überlassen die übrigen Aufgaben ergänzenden Paketen
\cite{comparative_spa_frameworks}.

Für die Auswahl bedeutet dieser Unterschied eine Abwägung. Ein
umfassendes Framework gibt eine einheitliche Struktur vor und
verringert den Entscheidungsaufwand, bindet das Projekt aber an die
Weiterentwicklung eines einzelnen Anbieters. Eine Bibliothek lässt
Freiheit bei der Zusammenstellung des Technologiestacks, verlagert
die Verantwortung für dessen Stimmigkeit jedoch in das Projekt
\cite{repolust_frontend_frameworks}. Die Aufgabenbereiche, die dabei
gesondert zu entscheiden sind, werden in den folgenden Abschnitten
einzeln betrachtet: die Wahl des Frameworks selbst
(Abschnitt~\ref{subsec:react}), die Verwaltung des
Anwendungszustands (Abschnitt~\ref{subsec:state-ansaetze}), die
Kommunikation mit dem Backend (Abschnitt~\ref{subsec:rest-api}), die
Gestaltung der Oberfläche (Abschnitt~\ref{subsec:ui-bibliotheken})
sowie die Verarbeitung von Formularen
(Abschnitt~\ref{subsec:formulare-validierung}).


% ---------------------------------------------------------------------
\subsection{Gegenüberstellung von React, Angular und Vue.js}
\label{subsec:react}

React ist eine JavaScript-Bibliothek zur Entwicklung
komponentenbasierter Benutzeroberflächen, die von Meta entwickelt und
gepflegt wird. Sie basiert auf dem Konzept des Virtual \acs{DOM},
wodurch Änderungen im Browser gezielt und performant aktualisiert
werden können, ohne die gesamte Seite neu zu laden. React hat sich als
eines der meistgenutzten Frontend-Frameworks etabliert
\cite{comparative_spa_frameworks}.

Angular verfolgt demgegenüber einen stark vorgebenden Ansatz mit
strikten Konventionen und liefert Routing, Formularverarbeitung und
\acs{HTTP}-Kommunikation bereits mit aus. React deckt lediglich die
View-Schicht ab und überlässt die Wahl der ergänzenden Bibliotheken
den Entwicklern. Vue.js liegt konzeptuell zwischen beiden Ansätzen,
verfügt jedoch über ein kleineres Ökosystem und eine geringere
Verbreitung in komplexen Unternehmensanwendungen
\cite{comparative_spa_frameworks}.

Eine systematische Gegenüberstellung der drei Lösungen anhand einer
Nutzwertanalyse über neun Kriterien, darunter \acs{API}-Stabilität,
Community-Unterstützung, Leistung und Testbarkeit, kommt zu dem
Ergebnis, dass React mit 75,15 Prozent die höchste Gesamtbewertung
erreicht, gefolgt von Vue.js mit 60,14 Prozent und Angular mit
57,30 Prozent. React wird dabei insbesondere die Stabilität seiner
Schnittstellen und die Breite seines Ökosystems zugerechnet, Vue.js
eine hohe Leistung bei Schwächen in der \acs{API}-Stabilität und
Angular eine klare Versionierungsstrategie bei eingeschränkter
interner Abwärtskompatibilität \cite{repolust_frontend_frameworks}.

Von den genannten Lösungen abzugrenzen sind Meta-Frameworks wie
Next.js, die React um serverseitiges Rendering und um eine
vorgegebene Projektstruktur ergänzen. Sie bieten Vorteile,
wenn Inhalte für Suchmaschinen zugänglich sein müssen oder eine
serverseitige Vorausführung erforderlich ist, erhöhen jedoch die
Komplexität der Betriebsumgebung \cite{pati_react_nextjs_2025}.
Tabelle~\ref{tab:frameworks} fasst die Gegenüberstellung zusammen.

\begin{table}[H]
\centering
\caption{Gegenüberstellung von React, Angular und Vue.js}
\label{tab:frameworks}
\small
\begin{tabular}{|p{2.6cm}|p{3.2cm}|p{3.2cm}|p{3.2cm}|}
\hline
\textbf{Kriterium} & \textbf{React} & \textbf{Angular} &
\textbf{Vue.js} \\
\hline
Einordnung & Bibliothek für die View-Schicht &
Umfassendes Framework & Framework mit modularem Kern \\
\hline
Funktions\-umfang & Ergänzende Bibliotheken erforderlich &
Routing, Formulare, \acs{HTTP} enthalten &
Kernfunktionen enthalten, Erweiterungen optional \\
\hline
Struktur\-vorgabe & Gering, hohe Gestaltungsfreiheit &
Hoch, strikte Konventionen & Mittel \\
\hline
Ökosystem & Sehr groß & Groß & Kleiner \\
\hline
Nutzwert nach \cite{repolust_frontend_frameworks} &
75,15 \% & 57,30 \% & 60,14 \% \\
\hline
\end{tabular}
\end{table}


% ---------------------------------------------------------------------
\subsection{Ansätze des State Managements}
\label{subsec:state-ansaetze}

Wie in Abschnitt~\ref{subsec:state-management} dargestellt, ist
zwischen clientseitigem Anwendungszustand und serverseitigem
Datenbestand zu unterscheiden. Für beide stehen unterschiedliche
Ansätze zur Verfügung.

Für den clientseitigen Zustand bietet React mit der Context
\acs{API} eine eingebaute Lösung, die Daten ohne Weiterreichen über
mehrere Komponentenebenen hinweg verfügbar macht. Sie eignet sich
für selten wechselnde Werte wie Sprache oder Farbschema, führt bei
häufigen Aktualisierungen jedoch zu unnötigen Neu-Renderings aller
eingebundenen Komponenten \cite{le_state_management}.

Redux verfolgt einen zentralisierten Ansatz mit unidirektionalem
Datenfluss: Zustandsänderungen werden ausschließlich über Aktionen
ausgelöst und in Reducern verarbeitet, wodurch jede Änderung
nachvollziehbar bleibt. Diese Nachvollziehbarkeit erkauft der Ansatz
mit einem hohen Anteil an wiederkehrendem Strukturcode und einer
Einbindung der Anwendung in einen Provider
\cite{redux_docs, le_state_management}.

Zustand verfolgt einen hook-basierten Ansatz, bei dem ein Store
unmittelbar als Hook definiert und ohne Provider-Wrapper in jeder
Komponente verwendet werden kann. Komponenten werden dabei nur dann
neu gerendert, wenn sich derjenige Teil des Zustands ändert, den sie
tatsächlich lesen \cite{zustand_docs}.

Für den serverseitigen Datenbestand besteht die Alternative zwischen
einer eigenen Verwaltung im clientseitigen Zustand und einer darauf
spezialisierten Bibliothek. Bei eigener Verwaltung müssen
Ladezustand, Fehlerfall und Aktualität für jede Anfrage einzeln
abgebildet werden. TanStack Query übernimmt diese Aufgaben
stattdessen zentral: Anfragen werden über einen Schlüssel
identifiziert, die Ergebnisse zwischengespeichert und der Zustand
jeder Verbindung den aufrufenden Komponenten einheitlich
bereitgestellt \cite{tanstack_query_docs}.
Tabelle~\ref{tab:state-management} stellt die Ansätze für den
clientseitigen Zustand gegenüber.

\begin{table}[H]
\centering
\caption{Gegenüberstellung der Ansätze für clientseitiges State
Management}
\label{tab:state-management}
\small
\begin{tabular}{|p{2.6cm}|p{3.2cm}|p{3.2cm}|p{3.2cm}|}
\hline
\textbf{Kriterium} & \textbf{Context \acs{API}} &
\textbf{Redux} & \textbf{Zustand} \\
\hline
Herkunft & In React enthalten & Externe Bibliothek &
Externe Bibliothek \\
\hline
Struktur\-code & Gering & Hoch (Aktionen, Reducer) & Gering \\
\hline
Provider erforderlich & Ja & Ja & Nein \\
\hline
Neu-Rendering & Alle eingebundenen Komponenten &
Selektiv über Selektoren & Selektiv je gelesenem Teilzustand \\
\hline
Geeignet für & Selten wechselnde Werte &
Komplexe, nachvollziehbare Zustandsflüsse &
Überschaubare Anzahl fachlich getrennter Stores \\
\hline
\end{tabular}
\end{table}


% ---------------------------------------------------------------------
\subsection{Kommunikationsansätze zwischen Frontend und Backend}
\label{subsec:rest-api}

\acfirst{REST} ist ein Architekturstil für Webschnittstellen, der auf
zustandslosen Anfragen, standardisierten \acs{HTTP}-Methoden und
ressourcenorientierter Adressierung basiert \cite{fielding_rest}.
Jede Anfrage enthält alle zur Verarbeitung notwendigen Informationen,
ohne dass der Server einen clientseitigen Sitzungszustand speichern
muss. Die Struktur der Antwort wird dabei serverseitig festgelegt.
Benötigt eine Ansicht Daten aus mehreren Ressourcen, sind mehrere
Anfragen erforderlich; benötigt sie nur einen Teil der gelieferten
Felder, werden dennoch alle übertragen.

GraphQL setzt an dieser Stelle an und verlagert die Bestimmung der
Antwortstruktur auf den Client. Über einen einzelnen Endpunkt
formuliert dieser eine Abfrage, die genau die benötigten Felder
benennt, auch über mehrere zusammenhängende Ressourcen hinweg.
Messungen zeigen, dass sich der Umfang der übertragenen Daten dadurch
erheblich reduzieren lässt und der Vorteil insbesondere bei hoher
Netzwerklatenz zum Tragen kommt. Demgegenüber weist \acs{REST}
Vorteile bei Antwortzeit und Speicherbedarf auf, insbesondere unter
hoher Last und bei häufig wiederholten, gleichartigen Anfragen; bei
komplexen Abfragen steigen unter GraphQL sowohl Antwortzeiten als
auch Fehlerraten \cite{diva_graphql_rest, lloyd_graphql_rest}.

Beiden Ansätzen gemeinsam ist, dass die Kommunikation vom Client
ausgeht. Muss der Server seinerseits Daten übermitteln, ohne dass
eine Anfrage vorliegt, etwa bei Benachrichtigungen oder Nachrichten
in Echtzeit, ist ein anfragebasiertes Verfahren ungeeignet.
WebSockets stellen hierfür eine dauerhafte, bidirektionale Verbindung
zwischen Client und Server bereit, über die beide Seiten jederzeit
Daten senden können, ohne dass der Client wiederholt Anfragen stellen
muss. WebSockets treten damit nicht an die Stelle von \acs{REST} oder
GraphQL, sondern ergänzen diese für Anwendungsfälle mit
Echtzeitanforderung.

Unabhängig vom gewählten Ansatz erfolgt der Zugriff im Browser über
einen \acs{HTTP}-Client. Axios ist ein verbreiteter,
Promise-basierter \acs{HTTP}-Client für JavaScript-Anwendungen. Er
erlaubt es, eine zentrale Instanz mit einer festen Basis-Adresse und
einheitlichen Standardeinstellungen zu konfigurieren, und stellt
sogenannte Interceptors bereit, mit denen ausgehende Anfragen und
eingehende Antworten an einer zentralen Stelle verarbeitet werden
können, etwa um jeder Anfrage automatisch einen
Authentifizierungstoken hinzuzufügen \cite{axios_docs}. Die im
Browser eingebaute Fetch-\acs{API} bietet vergleichbare
Grundfunktionen, jedoch ohne zentrale Konfiguration und ohne
Interceptors, sodass diese Aufgaben bei ihrer Verwendung selbst
umgesetzt werden müssen. Tabelle~\ref{tab:kommunikation} fasst die
Kommunikationsansätze zusammen.

\begin{table}[H]
\centering
\caption{Gegenüberstellung der Kommunikationsansätze}
\label{tab:kommunikation}
\small
\begin{tabular}{|p{2.6cm}|p{3.2cm}|p{3.2cm}|p{3.2cm}|}
\hline
\textbf{Kriterium} & \textbf{\acs{REST}} & \textbf{GraphQL} &
\textbf{WebSockets} \\
\hline
Initiative & Client & Client & Client und Server \\
\hline
Endpunkte & Mehrere, ressourcen\-orientiert & Einer &
Kanal\-basiert \\
\hline
Antwort\-struktur & Serverseitig festgelegt &
Clientseitig bestimmt & Anwendungs\-spezifisch \\
\hline
Stärke & Antwortzeit, Zwischen\-speicherung &
Geringes Daten\-volumen, eine Anfrage je Ansicht &
Übertragung in Echtzeit \\
\hline
Schwäche & Über\-tragung nicht benötigter Felder &
Höhere Antwortzeit bei komplexen Abfragen &
Dauerhafte Verbindung erforderlich \\
\hline
\end{tabular}
\end{table}


% ---------------------------------------------------------------------
\subsection{Styling-Ansätze und UI-Bibliotheken}
\label{subsec:ui-bibliotheken}

Für die Gestaltung der Oberfläche lassen sich drei Ansätze
unterscheiden. Komponentenbasierte \acs{CSS}-Frameworks wie Bootstrap
liefern fertig gestaltete Elemente aus, die unmittelbar verwendbar
sind, deren Anpassung an ein eigenes Erscheinungsbild jedoch Aufwand
erfordert. Bei \acs{CSS}-in-JS werden Stile innerhalb des
JavaScript-Codes einer Komponente definiert, wodurch Stil und Logik
zusammenliegen, zur Laufzeit jedoch zusätzlicher Verarbeitungsaufwand
entsteht.

Utility-first-Frameworks wie Tailwind CSS verfolgen einen dritten
Ansatz: Sie definieren keine fertigen Komponenten, sondern atomare
Stilklassen für Abstände, Farben und Typografie, die direkt im
\acs{HTML}-Markup kombiniert werden. Dadurch entfällt das Schreiben
separaten \acs{CSS}, und Stile sind unmittelbar im Markup sichtbar
\cite{tailwindcss_docs}. Vergleichende Untersuchungen ordnen
utility-first-Frameworks eine höhere Kontrolle und Effizienz zu,
weisen zugleich aber auf eine steilere Lernkurve und längere
Klassenangaben im Markup hin \cite{utility_first_comparison}.

Ergänzend zu diesen Ansätzen existieren Komponentensammlungen wie
shadcn/ui, die zugängliche, anpassbare React-Komponenten auf Basis
von Radix UI-Primitives bereitstellen und mit Tailwind CSS gestaltet
sind. Im Unterschied zu klassischen Komponentenbibliotheken wird
shadcn/ui nicht als Paketabhängigkeit installiert; die
Komponentenquellcodes werden stattdessen direkt in das Projekt
kopiert, wodurch die vollständige Kontrolle über den Code beim
Projekt verbleibt \cite{shadcn_docs}.


% ---------------------------------------------------------------------
\subsection{Formularverarbeitung und Validierung}
\label{subsec:formulare-validierung}

Formulare bilden in E-Commerce-Anwendungen einen zentralen
Interaktionspunkt, etwa bei Registrierung und Anmeldung oder der
Verwaltung von Produkten und Dienstleistungen. React Hook Form
ist eine Bibliothek zur Formularverwaltung in React-Anwendungen,
die Eingabefelder über Hooks registriert und den Formularzustand
mit möglichst wenigen Neu-Renderings verwaltet
\cite{react_hook_form_docs}. Zod ist eine TypeScript-orientierte
Bibliothek zur Schema-Validierung, mit der Datenstrukturen
deklarativ definiert und zur Laufzeit geprüft werden; aus den
Schemata werden zugleich statische Typen abgeleitet
\cite{zod_docs}. Beide Bibliotheken lassen sich über einen
sogenannten Resolver koppeln, sodass Formulareingaben
clientseitig gegen ein Zod-Schema validiert werden, bevor eine
Anfrage an das Backend gestellt wird \cite{react_hook_form_docs}.


% =====================================================================
% HERKUNFT DES TEXTES -- bitte lesen
%
% AUS DEINEM BISHERIGEN TEXT UEBERNOMMEN (unveraendert oder nur
% entschaerft, weil die Bazario-Bezuege raus mussten):
%   - 2.3.2, Absatz 1 und 2 (React, Virtual DOM, Angular/Vue)
%   - 2.3.4, REST-Definition nach Fielding
%   - 2.3.4, Beschreibung von Axios und Interceptors
%   - 2.3.5, Tailwind CSS und shadcn/ui
%   - 2.3.6 vollstaendig, unveraendert
%
% VON MIR NEU FORMULIERT (Literaturbasis vorhanden, aber die
% Formulierung und die Zuspitzung stammen nicht von dir):
%   - 2.3.1 vollstaendig
%   - 2.3.2, Absatz 3 (Nutzwertanalyse) und Absatz 4 (Next.js)
%   - 2.3.3 vollstaendig
%   - 2.3.4, Absaetze zu GraphQL und WebSockets
%   - 2.3.5, Absatz 1 (Bootstrap und CSS-in-JS)
%   - alle drei Tabellen
%
% Nach deiner Regel musst du diese Passagen selbst nachziehen. Sie
% sind belegt, aber es sind nicht deine Saetze.
%
%
% TODO (Roudy):
%
% 1. NEUE ABKUERZUNG: \acfirst{DOM} in 2.3.1 setzt voraus, dass DOM
%    in Kapitel/0_Abkuerzungsverzeichnis definiert ist. Pruefen und
%    ggf. anlegen. Im bisherigen Text kam "Virtual DOM" nur
%    ausgeschrieben vor.
%
% 2. Die Prozentwerte in 2.3.2 und in der Tabelle (75,15 / 60,14 /
%    57,30) stammen aus der Masterarbeit von Repolust. Pruefe im PDF,
%    ob du sie so uebernehmen willst. Wichtig: Die Arbeit bewertet
%    Eignung fuer LEGACY-Projekte anhand von Abwaertskompatibilitaet.
%    Das ist nicht exakt dein Anwendungsfall. Entweder du erwaehnst
%    den Untersuchungsrahmen im Satz, oder du laesst die Zahlen weg
%    und nutzt die Arbeit nur qualitativ. Ich wuerde den Rahmen
%    nennen, sonst ist es angreifbar.
%
% 3. Die GraphQL-Aussagen in 2.3.4 stuetzen sich auf zwei Quellen,
%    deren Autor und Jahr in literatur_neu_fuer_2_3.bib noch als TODO
%    stehen. Vor der Abgabe ergaenzen.
%
% 4. Zu 2.3.3: Ich habe den IJSAT-Artikel bewusst NICHT zitiert
%    (Begruendung in literatur_neu_fuer_2_3.bib). Der Abschnitt
%    stuetzt sich auf le_state_management, redux_docs, zustand_docs
%    und tanstack_query_docs. Wenn dir das zu duenn erscheint, sag
%    Bescheid.
%
% 5. In 2.3.5 habe ich CSS-in-JS ohne eigene Quelle beschrieben, weil
%    ich dafuer keine belastbare gefunden habe. Der Absatz ist
%    inhaltlich unstrittig, aber unbelegt. Entweder Quelle nachtragen
%    oder den CSS-in-JS-Satz streichen und nur Bootstrap gegen
%    utility-first stellen.
%
% 6. Pruefen, ob \acs{REST} beim ersten Vorkommen in 2.3.4 wirklich
%    \acfirst sein muss. In deinem alten Text stand \acfirst{REST}
%    dort. Da REST aber schon in 2.1 und 2.2 vorkommt, kann die
%    Langform bereits verbraucht sein.
% =====================================================================

\section{Qualitätssicherung in der Frontend-Entwicklung}
\label{sec:qualitaetssicherung}

\subsection{Teststrategie: Unit-, Integrations- und E2E-Tests}
\label{subsec:teststrategie}

Eine strukturierte Teststrategie ist ein wesentlicher
Bestandteil der Qualitätssicherung in der
Softwareentwicklung. Agile Entwicklungsansätze empfehlen
eine Kombination aus verschiedenen Testebenen, häufig
als Testpyramide bezeichnet, um sowohl einzelne
Komponenten als auch das Zusammenspiel des Gesamtsystems
abzusichern \cite{cohn_agile}. Die Testpyramide umfasst
an der Basis zahlreiche schnelle Unit-Tests, darüber
Integrationstests und an der Spitze eine begrenzte
Anzahl End-to-End-Tests.

Unit-Tests und Integrationstests werden mit Vitest
durchgeführt. Vitest ist ein nativer Test-Runner für
Vite-basierte React-Anwendungen und ist vollständig
kompatibel mit dem Vite-Build-System ohne zusätzliche
Konfiguration. Während Unit-Tests einzelne Komponenten
und Funktionen isoliert prüfen, testen Integrationstests
das Zusammenspiel mehrerer Komponenten, etwa die
Interaktion eines Formulars mit dem globalen State
\cite{vitest_docs}.

End-to-End-Tests simulieren reale Nutzerinteraktionen
im Browser und prüfen vollständige Anwendungsabläufe.
Playwright ermöglicht die automatisierte Ausführung
von Nutzerflows in mehreren Browsern gleichzeitig.
Für \textit{Bazario} mit seinen verschiedenen
Nutzerrollen bietet Playwright entscheidende Vorteile,
da rollenspezifische Abläufe wie Buchung, Checkout
und Adminfreigabe browserübergreifend getestet werden
können \cite{playwright_docs}.

\subsection{Usability und User Experience}
\label{subsec:usability}

Neben funktionaler Korrektheit ist die
Gebrauchstauglichkeit einer Webanwendung ein zentrales
Qualitätsmerkmal. Heuristische Evaluation ist eine
etablierte Methode zur Bewertung der Usability, bei der
Experten die Benutzeroberfläche anhand definierter
Heuristiken auf Usability-Probleme untersuchen. Im
Kontext von E-Commerce-Plattformen wurden spezifische
Heuristiken entwickelt, die auf die besonderen
Anforderungen digitaler Marktplätze abgestimmt sind,
darunter Kriterien wie Navigationskonsistenz,
Transparenz von Preisen und Prozessen sowie die
Effizienz des Checkout-Vorgangs
\cite{abdulaziz_heuristic_marketplace}.

User Experience umfasst dabei neben der reinen
Gebrauchstauglichkeit auch emotionale und
wahrnehmungsbezogene Aspekte der Nutzung. Für
E-Commerce-Websites wurden \acs{UX}-spezifische Heuristiken
erarbeitet, die Aspekte wie Vertrauen, Orientierung,
Effizienz des Checkout-Prozesses und visuelle Konsistenz
berücksichtigen \cite{bonastre_ux_heuristics}. Diese
Erkenntnisse fließen in die Konzeption und Bewertung
des Frontends von \textit{Bazario} ein und bilden die
Grundlage für die Usability-Evaluation in
Kapitel~\ref{chap:testen}.